\documentclass[11pt]{article}
\usepackage[a4paper,margin=25mm]{geometry}
\usepackage{amsmath,amssymb,amsthm,hyperref}
\newtheorem{lemma}{Lemma}
\newtheorem{corollary}{Corollary}
\title{All-order localized derivative bounds for the Laguerre Poisson kernel}
\author{Auxiliary analytic argument; independent mathematical reviews completed}
\date{30 September 2026}
\begin{document}\maketitle
The full argument passed separate mathematical reviews by the size-bounds
and rational-design collaborators. This is not formal verification.
This note establishes an analytic input for a possible improvement of
endpoint-profile bounds. It does not establish a new endpoint-profile
rate: polynomial truncation, integration over the row measure, and the
final monotonicity argument remain separate tasks.

For $0<\tau<1$ define the positive Laguerre Poisson kernel
\[
 K_\tau(x,a)=\frac1\tau
  \exp\!\left[-\frac{(1-\tau)(x+a)}\tau\right]
  I_0\!\left(\frac{2\sqrt{(1-\tau)xa}}\tau\right).
\]
The Hille--Hardy identity, at parameter $\alpha=0$, says
\[
 K_\tau(x,a)=\sum_{j=0}^{\infty}(1-\tau)^j L_j(x)L_j(a).
\]
See \href{https://dlmf.nist.gov/18.18.E27}{DLMF 18.18.27}.
Although written using a square root, $K_\tau(\cdot,a)$ is entire,
as is immediate from the power series of $I_0$.

\begin{lemma}\label{derivatives}
Fix $0<\delta<L<\infty$. There are constants $C,c,\tau_0>0$,
depending only on this interval, such that for $x,a\in[\delta,L]$,
$0<\tau<\tau_0$, and every integer $k\ge1$,
\[
 \frac{|\partial_x^k K_\tau(x,a)|}{k!}
 \le C K_{4\tau}(x,a)
 \left[\left(\frac C{\sqrt{k\tau}}\right)^k+C^k\right].       \tag{1}
\]
Moreover, whenever $|z|\le c$,
\[
 |K_\tau(x+z,a)|\le C K_{4\tau}(x,a)
                              \exp(C|z|^2/\tau).             \tag{2}
\]
\end{lemma}
\begin{proof}
Consider first the radial Gaussian kernel
\[
 B_\tau(y,a)=\tau^{-1}e^{-(y+a)/\tau}
                             I_0(2\sqrt{ya}/\tau).
\]
For real positive arguments in any fixed compact interval,
the elementary integral formula
\[
 I_0(t)=\frac1{2\pi}\int_{-\pi}^{\pi}e^{t\cos\phi}\,d\phi
\]
and $1-\cos\phi\asymp\phi^2$ on $[-\pi,\pi]$ show that
\[
 I_0(t)\asymp e^t/\sqrt t\quad(t\ge1),\qquad
 B_\tau(y,a)\asymp\tau^{-1/2}
                   e^{-(\sqrt y-\sqrt a)^2/\tau}.          \tag{3}
\]
The constants in the second comparison are uniform on that compact
interval for sufficiently small $\tau$.

Let $y>0$, $|z|\le c$, and write $\sqrt{y+z}=u+iv$, choosing the
branch with $u>0$. On the relevant fixed compact intervals, choosing
$c$ small gives
\[
 |u-\sqrt y|+|v|\le C|z|.
\]
Taking absolute values inside the integral formula for $B_\tau$ gives
the exact useful domination
\[
 |B_\tau(y+z,a)|\le e^{v^2/\tau}B_\tau(u^2,a).
\]
Since
$(u-\sqrt a)^2\ge\tfrac12(\sqrt y-\sqrt a)^2-C|z|^2$,
formula (3) therefore gives
\[
 |B_\tau(y+z,a)|\le C\tau^{-1/2}
 \exp\!\left[-\frac{(\sqrt y-\sqrt a)^2}{2\tau}
                       +\frac{C|z|^2}\tau\right].         \tag{4}
\]

Now $K_\tau(x,a)=e^a B_\tau((1-\tau)x,a)$.
Write $y_1=(1-\tau)x$ and $y_4=(1-4\tau)x$.
They stay in a fixed compact positive interval, and
$|\sqrt{y_1}-\sqrt{y_4}|\le C\tau$. Thus
\[
 (\sqrt{y_1}-\sqrt a)^2
 \ge\tfrac12(\sqrt{y_4}-\sqrt a)^2-C\tau^2.
\]
Use this in (4), then the lower bound (3) for $B_{4\tau}(y_4,a)$.
The additional factor $e^{C\tau}$ is bounded. This proves (2).

Cauchy's formula on a circle of radius $R\le c$ consequently yields
\[
 \frac{|\partial_x^kK_\tau(x,a)|}{k!}
 \le C K_{4\tau}(x,a)R^{-k}e^{CR^2/\tau}.
\]
Choose $R=\min(c,\sqrt{k\tau/(2C)})$.
If the second argument is smaller, the result is bounded by
$C K_{4\tau}(C/\sqrt{k\tau})^k$.
Otherwise $k\tau\ge2Cc^2$, so $e^{Cc^2/\tau}\le e^{k/2}$,
and the result is at most $C K_{4\tau}C^k$.
This proves (1).
\end{proof}

\begin{corollary}[Centered polarization error on a fixed interval]
Let $z_1,\ldots,z_N\in[\delta,L]$ have mean $v$ and range at most $H$.
Set $a_k=e_k(z_1-v,\ldots,z_N-v)/\binom Nk$.
For any integer $d\le N$, if
\[
 H\le c,\qquad H/\sqrt{N\tau}\le c,
\]
then, uniformly in $a\in[\delta,L]$,
\[
 \left|\sum_{k=2}^d
   \frac{\partial_v^kK_\tau(v,a)}{k!}a_k\right|
 \le C K_{4\tau}(v,a)H^2
                      \left(\frac1{N\tau}+\frac1N\right). \tag{5}
\]
Constants depend only on $\delta,L$.
\end{corollary}
\begin{proof}
The elementary centered coefficient bound gives
$|a_k|\le(CH\sqrt{k/N})^k$; a proof is in
\path{independent_directions/marked_cancellation_identity.tex}.
For the first term in (1), sum the geometric series in
$CH/\sqrt{N\tau}$ beginning at $k=2$.
For the second term, use $k\le N$ to write
\[
 (CH)^k(k/N)^{k/2}\le (CH)^k k/N,
\]
and sum. This gives (5).
\end{proof}

In (5) the expression is a finite differential sum, not a claim that
the infinite Poisson kernel is a polynomial. Applying it to exact
finite moment identities requires a justified polynomial truncation.
Likewise, the constants here are not uniform as $\delta\downarrow0$.
\end{document}
